\section{Проектирование и разработка системы управления молочным животноводством}

\subsection{Архитектура системы}

Разработанная информационная система имеет трёхзвенную архитектуру, состоящую из серверной части (backend), клиентского веб-приложения (frontend) и модуля машинного обучения (analytics-ml). Компоненты развёрнуты в контейнерах Docker и объединены через обратный прокси Nginx.

На рисунке~\ref{arch} представлена архитектура системы. Nginx является единой точкой входа, принимающей HTTP-запросы от клиентов. Запросы к статическим ресурсам и страницам маршрутизируются к контейнеру frontend (SvelteKit), запросы к API --- к контейнеру backend (Rust/Axum). Backend обращается к PostgreSQL для хранения данных, к Redis для кэширования и rate limiting, и к analytics-ml для получения прогнозов машинного обучения. Синхронизация с облачным API Lely выполняется периодически по таймеру внутри backend.

\begin{figure}
  \centering
  \includegraphics[width=0.95\textwidth,height=0.85\textheight,keepaspectratio]{container.png}
  \caption{Архитектура информационной системы}
  \label{arch}
\end{figure}

Архитектурные решения обоснованы несколькими соображениями. Разделение ответственности обеспечивает чёткое разграничение: backend отвечает за бизнес-логику, frontend --- за пользовательский интерфейс, analytics-ml --- за модели машинного обучения. Независимое масштабирование позволяет запускать несколько экземпляров backend за балансировщиком нагрузки без изменения остальных компонентов. Технологическая гетерогенность даёт возможность использовать наиболее подходящий стек для каждого компонента (Rust для серверной части, Svelte для клиента, Python для ML). Изоляция сбоев, обеспечиваемая контейнерной изоляцией Docker, гарантирует, что сбой одного компонента не приводит к отказу всей системы.

\subsection{Модель данных}

Архитектура системы определяет структуру хранилища данных, к описанию которой мы переходим. База данных PostgreSQL содержит более 35 таблиц, обеспечивающих хранение всех сущностей предметной области. Перечень таблиц с назначением приведён в приложении~\ref{app-db}, полная ER-модель --- в приложении~\ref{app-er}. На рисунке~\ref{er} представлена модель данных системы.

\begin{figure}
  \centering
  \includegraphics[width=0.95\textwidth,height=0.85\textheight,keepaspectratio]{database-model.png}
  \caption{Модель данных информационной системы}
  \label{er}
\end{figure}

Основные таблицы базы данных и их назначение:

\textbf{users} --- пользователи системы. Содержит имя пользователя, хеш пароля (bcrypt), роль (admin/user), флаг обязательной смены пароля и дату создания. Пароли хранятся в виде хешей с использованием алгоритма bcrypt с автоматической генерацией соли, что исключает утечку паролей при компрометации базы данных.

\textbf{animals} --- реестр животных. Содержит уникальный идентификатор, инвентарный номер (life\_number), кличку, пол, дату рождения, породу, местоположение, статус активности и множество других полей. Поддерживается флаг \texttt{active} для деактивации животных без удаления записей, что обеспечивает сохранение истории.

\textbf{milk\_day\_productions} --- суточные надои. Привязаны к животному и дате. Содержат объём молока (кг), жир (\%), белок (\%), лактозу (\%), количество соматических клеток (\acrshort{scc}), энергию (MJ) и жирово-белковое соотношение (\acrshort{fpr}). Композитный индекс (animal\_id, date) обеспечивает быстрый поиск надоев за период.

\textbf{milk\_visits} --- визиты на доение. Содержат информацию о каждом доении: время начала и окончания, объём молока, скорость доения (кг/мин), данные по четвертям вымени (объём, время, процент выдаивания), флаг успешности доения и идентификатор робота.

\textbf{milk\_quality} --- результаты лабораторных анализов молока. Включает количество соматических клеток (\acrshort{scc}), жир, белок, лактозу, мочевину, бактерии.

\textbf{calvings, calves, inseminations, pregnancies, heats, dry\_offs} --- данные о воспроизводстве. Каждая таблица связана с животным через внешний ключ. Поддерживается полная история всех событий воспроизводства, включая дату события, результат и примечания.

\textbf{feed\_types, feed\_groups, feed\_day\_amounts, feed\_visits} --- данные о кормлении. Типы кормов описывают рационы с указанием стоимости и единиц измерения, группы --- привязку кормов к группам животных, дневные объёмы --- плановое и фактическое потребление, визиты --- отдельные посещения кормостанции с временем и объёмом.

\textbf{activities, ruminations, grazing\_data} --- данные фитнес-мониторинга и пастьбы. Активность измеряется в шагах за период, жвачка --- в минутах жевания за период.

\textbf{alerts} --- система оповещений. Оповещения создаются автоматически на основе настраиваемых пороговых значений. Каждое оповещение содержит категорию (milk, health, reproduction, feed, equipment), степень важности (high, medium, low), статус (active, acknowledged, resolved), идентификатор животного и текстовое описание.

\textbf{lely\_sync\_state} --- состояние синхронизации с Lely. Хранит дату последней успешной синхронизации для каждого типа данных, что позволяет инкрементально обновлять информацию и минимизировать объём данных при каждой синхронизации.

\textbf{lely\_config} --- конфигурация интеграции с Lely. Хранит URL API, зашифрованные логин и пароль, ключ фермы и настройки синхронизации. Учётные данные шифруются с использованием AES-256-GCM~\cite{AES-GCM}.

Для эффективной работы с временными рядами (надои, активность, жвачка) используется расширение TimescaleDB~\cite{TimescaleDB}. Таблицы с временными рядами преобразованы в гипертаблицы с автоматической партицировкой по времени (по умолчанию --- 7 дней на чанк), что ускоряет запросы за определённый период, так как сканируются только релевантные чанки. Непрерывные агрегаты (continuous aggregates) обеспечивают предвычисление суточных и часовых агрегатов, обновляемых в фоновом режиме.

Индексация включает уникальные ограничения (инвентарный номер животного), композитные индексы (животное + дата) и частичные индексы (только активные животные). Расширение pg\_trgm~\cite{PgTrgm} обеспечивает полнотекстовый поиск по кличкам животных с использованием триграммного сходства.
